\section{Introdução}
\label{sec:introducao}
% Parágrafo introdutório contextualizando o tema. 
% Apresentar o problema de forma geral. ~3 parágrafos.

\subsection{Justificativa}
% Por que este trabalho é relevante? Qual o impacto prático/acadêmico?
% Mencionar demanda da área, lacuna na literatura ou necessidade social. ~2 parágrafos.

\subsection{Objetivos}

\subsubsection{Objetivo Geral}
% Uma frase clara descrevendo o objetivo principal do trabalho.

\subsubsection{Objetivos Específicos}
\begin{itemize}
  \item [Objetivo específico 1 — verbo no infinitivo];
  \item [Objetivo específico 2 — verbo no infinitivo];
  \item [Objetivo específico 3 — verbo no infinitivo];
  \item [Objetivo específico 4 — verbo no infinitivo].
\end{itemize}

\subsection{Estrutura do Trabalho}
% Descrever brevemente o que cada seção seguinte contém.
% Ex.: "A Seção II apresenta o referencial teórico..."
